\documentclass{article}
\usepackage[utf8]{inputenc}
\usepackage{hyperref}
\begin{document}
Within a PIGEAN factor model (cfde-inc-v2), the GOBP\_STEROL\_METABOLIC\_PROCESS-labeled Factor1 for coronary artery disease (CAD) was examined to assess whether LPL qualifies as a candidate gene relevant to the open question of whether CAD PGS×context interactions reflect upstream causal amplification of genetic risk or downstream biochemical readouts of incipient disease.\cite{dapper:KnowledgeGap.zNV20nhHamt-a4CeAktQQPoAivOJe6xk-r1} Among the four retained direct candidate genes on this factor—LPL, SIRT3, FURIN, and TIMP3—LPL carries the highest factor loading (0.3439) and combined PIGEAN score (5.9, log\_bf=3.9, prior=1.99), identifying it as the top candidate in the sterol-metabolism-associated component of the CAD genetic architecture in this model.\cite{dapper:Claim.oHAogBRdetx6fDiakfRSC0LunxjGDR0E-r1,dapper:Claim.NZs1bmu2UfabvcU9OGLiAWrLIt_XDE61-r1} This factor loading is a quantitative association and does not establish causal direction; the CFDE factor association data cannot determine whether LPL's prominence reflects upstream participation in lipid-related genetic risk amplification or a downstream metabolic readout of disease, and the knowledge gap remains open pending evidence such as Mendelian randomization or prospective biomarker data stratified by polygenic score.\cite{dapper:Claim.NZs1bmu2UfabvcU9OGLiAWrLIt_XDE61-r1}
\bibliographystyle{plain}
\bibliography{references}
\end{document}
